\section{\filetotitle}\label{sec:numerical_schemes_for_time_dependent_schrodinger_equation}
This section focuses on numerical solutions to \acrshort{tdse}. In the described methods, we work in a position representation $\Psi\left(\symbf{R},t\right):=\braket{\symbf{R}\vert\Psi\left(t\right)}$. More importantly, we expand the wavefunction $\ket{\Psi\left(t\right)}$ in a set of diabatic basis functions $\lbrace\tilde\phi_I\rbrace_{I=1}^N$ defined in Sec~\ref{sec:adiabatic_vs_diabatic_representation}, so that $\Psi\left(\symbf{R},t\right)$ is an N-component vector in a electronic diabatic subspace.
\subsection{Split-Operator Method}\label{sec:split_operator_method}
For a time-independent Hamiltonian, the solution to the \acrshort{tdse} can be written as
\begin{equation}\label{eq:tdse_propagator}
\ket{\Psi\left(t\right)}=e^{-i t H}\ket{\Psi\left(0\right)}.
\end{equation}
The action of the operator $e^{-i t H}$ on the initial wavefunction $\ket{\Psi\left(0\right)}$ propagates the state from time $0$ to time $t$, carrying all the information about the time evolution of the system. Since $H$ is self-adjoint, the spectral theorem yields a unitary propagator, ensuring $\lvert\Psi(t)\rvert=\lvert\Psi(0)\rvert$. This exponential operator contains the complete description of how the system evolves with time. In practical applications, however, evaluating $e^{-i t H}$ analytically is rarely straightforward for nontrivial Hamiltonians, which is exactly the problem the Split-Operator method solves.

For an N-state system, the wavefunction in the diabatic basis can be written as
\begin{equation}\label{eq:diabatic_wavefunction}
\Psi\left(\symbf{R},t\right)=\sum_{I=1}^N c_I\left(\symbf{R},t\right)\ket{\tilde\phi_I}
\end{equation}
where $c_I\left(\symbf{R},t\right)$ are the expansion coefficients of the wavefunction in the diabatic basis. The Hamiltonian operator $H$ in the diabatic basis takes the form
\begin{equation}\label{eq:diabatic_hamiltonian}
H=\tilde{T}_{\symrm{n}}+\tilde{V}=-\frac{1}{2m}\nabla_\mathbf{R}^2\otimes I_N+\tilde{V}
\end{equation}
where $I_N$ is the identity operator on the $N$-dimensional diabatic electronic space,
$\tilde{T}_{\symrm{n}}$ is the nuclear kinetic energy operator, and $\tilde{V}$ is the
potential energy operator represented by an $N\times N$ matrix in the diabatic electronic basis. We aim to reduce the action of the propagator on the diabatic wavefunction in Eq~\eqref{eq:diabatic_wavefunction} to a simple matrix multiplication. The main difficulty is the differential form of the kinetic operator in the equation~\eqref{eq:diabatic_hamiltonian}. To handle this, we use a Fourier-based method for applying the kinetic operator $\tilde{T}_{\symrm{n}}$ on a wavefunction $\Psi\left(\symbf{R},t\right)$. Taking the \acrfull{ft} of $\tilde{T}_{\symrm{n}}\Psi\left(\symbf{R},t\right)$ yields
\begin{equation}\label{eq:kinetic_fourier_method}
\mathcal{F}\left\lbrace \tilde{T}_{\symrm{n}}\Psi\left(\symbf{R},t\right)\right\rbrace=\mathcal{F}\left\lbrace-\frac{1}{2m}\frac{\partial^2}{\partial \symbf{R}^2}\Psi\left(\symbf{R},t\right)\right\rbrace=\frac{k^2}{2m}\mathcal{F}\left\lbrace\Psi\left(\symbf{R},t\right)\right\rbrace,
\end{equation}
where $k$ is the coordinate in the Fourier space. In other words, to apply the kinetic operator on the wavefunction, we need to \acrshort{ft} the wavefunction, multiply it by $\frac{k^2}{2m}$, and then apply the \acrfull{ift}. In Fourier space, $\tilde{T}_{\symrm{n}}$ becomes diagonal multiplication operator $\frac{k^2}{2m}$. For notational simplicity we denote both the coordinate-space operator and its Fourier-space representation by the same symbol. This method is computationally efficient, since the \acrshort{ft} can be done using the \acrfull{fft} algorithm, which has a complexity of $\mathcal{O}(N\log N)$, where $N$ is the number of grid points (contrary to the straightforward \acrfull{dft} with $\mathcal{O}(N^2)$ complexity).

A remaining challenge is the non-commutativity of the potential and kinetic operators. Because $\tilde{V}$ and $\tilde{T}_{\symrm{n}}$ do not commute we cannot factorize the propagator as
\begin{equation}
e^{-i t H}\neq e^{-i t\tilde{V}}e^{-i t\tilde{T}_{\symrm{n}}}.
\end{equation}
If such a factorization were possible, we could apply the exponential of $\tilde{V}$ in the position space and the exponential of $\tilde{T}_{\symrm{n}}$ momentum space separately. To address this, we split the total propagation time $t$ into $n$ short steps $\Delta t$, so the full propagator becomes a product of short-time propagators as
\begin{equation}
U(t)=e^{-i t H}=\left(e^{-i\Delta t H}\right)^n,\qquad t=n\Delta t
\end{equation}
where $\Delta t$ is a fixed time step. Now, each individual short-time propagator acts on the wave function and causes a small evolution. We can then use the symmetric splitting approximation
\begin{equation}
U(\Delta t)=e^{-i\frac{\Delta t}{2}\tilde{V}}e^{-i\Delta t\tilde{T}_{\symrm{n}}}e^{-i\frac{\Delta t}{2}\tilde{V}}+\mathcal{O}(\Delta t^3),
\end{equation}
which is valid for small $\Delta t$. Now we have everything we need to propagate the wavefunction in time. The formula for propagating the wavefunction from time $t$ to time $t+\Delta t$ is
\begin{equation}
\Psi\left(\symbf{R},t+\Delta t\right)=e^{-i\frac{\Delta t}{2}\tilde V}\mathcal{F}^{-1}\left\lbrace e^{-i\Delta t\tilde{T}_{\symrm{n}}}\mathcal{F}\left\lbrace e^{-i\frac{\Delta t}{2}\tilde V}\Psi\left(\symbf{R},t\right)\right\rbrace\right\rbrace,
\end{equation}
where $\tilde{T}_{\symrm{n}}$ is the kinetic operator in the momentum space, and $\tilde{V}$ is the potential operator in the position space. In a discrete grid representation these become matrices acting on the sampled wavefunction.
\subsection{Crank--Nicolson Scheme}\label{sec:crank_nicolson_scheme}
The Split-Operator method in Sec.~\ref{sec:split_operator_method} exploits the simple action
of the kinetic and potential operators in Fourier and position representation, respectively,
and approximates the short-time propagator by a symmetric factorization of
$e^{-i\Delta t H}$. An alternative widely used time-integration scheme is the
Crank--Nicolson method, which can be viewed as applying a rational (Padé) approximation to
the propagator.

We again consider the \acrshort{tdse} for a time-independent, self-adjoint Hamiltonian $H$,
\begin{equation}\label{eq:tdse_cn}
i\frac{\symrm{d}}{\symrm{d}t}\ket{\Psi(t)}=H\ket{\Psi(t)},
\end{equation}
acting on the separable Hilbert space $\mathcal{H}$ introduced in Sec.~\ref{sec:hilbert_spaces_and_quantum_states}. Let $t_n=n\Delta t$ denote a uniform time grid with step $\Delta t>0$ and define
\begin{equation}
\ket{\Psi^n}:=\ket{\Psi(t_n)}.
\end{equation}
The Crank–Nicolson scheme is obtained by approximating the time derivative by a centered second-order accurate finite difference at the midpoint $t_{n+\frac12}=t_n+\Delta t/2$ and evaluating the right-hand side via the trapezoidal rule (midpoint average) as
\begin{equation}\label{eq:cn_trapezoidal}
\frac{\ket{\Psi^{n+1}}-\ket{\Psi^n}}{\Delta t}=-\frac{i}{2}\left(H\ket{\Psi^{n+1}}+H\ket{\Psi^n}\right).
\end{equation}
Both sides of Eq.~\eqref{eq:cn_trapezoidal} approximate the exact quantities at the midpoint
with an $\mathcal{O}(\Delta t^2)$ error. Rearranging Eq.~\eqref{eq:cn_trapezoidal} yields the linear system
\begin{equation}\label{eq:cn_linear_system}
\left(I+\frac{i\Delta t}{2}H\right)\ket{\Psi^{n+1}}=\left(I-\frac{i\Delta t}{2}H\right)\ket{\Psi^n},
\end{equation}
where $I$ is the identity operator on $\mathcal{H}$. For each time step one solves Eq.~\eqref{eq:cn_linear_system} for $\ket{\Psi^{n+1}}$ given $\ket{\Psi^n}$.

Formally, Eq.~\eqref{eq:cn_linear_system} defines a one-step propagation
\begin{equation}\label{eq:cn_propagator_operator}
\ket{\Psi^{n+1}}=U_{\symrm{CN}}(\Delta t)\ket{\Psi^n},\qquad U_{\symrm{CN}}(\Delta t)=\left(I+\frac{i\Delta t}{2}H\right)^{-1}\left(I-\frac{i\Delta t}{2}H\right).
\end{equation}
The operator $U_{\symrm{CN}}(\Delta t)$ is equivalent to the $[1,1]$ Padé approximation to the exact propagator $e^{-i\Delta t H}$:
\begin{equation}
U_{\symrm{CN}}(\Delta t)=e^{-i\Delta t H}+\mathcal{O}(\Delta t^3).
\end{equation}
Because $H$ is self-adjoint, $U_{\symrm{CN}}(\Delta t)$ is unitary so that $\lvert\Psi^{n+1}\rvert=\lvert\Psi^n\rvert$ and the total probability is conserved, as required by the unitary time evolution discussed in Sec.~\ref{sec:spectral_theory_in_hilbert_spaces}.
\subsection{Imaginary Time Propagation}\label{sec:imaginary_time_propagation}
Imaginary time propagation is a numerical technique for obtaining the ground state of a
quantum system by evolving it in imaginary time instead of real time. By replacing the real
time variable by an imaginary one, contributions from higher-energy eigenstates are
suppressed exponentially faster than that of the ground state, so that repeated
propagation and renormalization project an initial state onto the ground state. As with real-time propagation, its practical applicability is restricted by the cost of repeatedly applying the Hamiltonian to the wavefunction.

We again consider the \acrshort{tdse} for a time-independent, self-adjoint Hamiltonian
$H$ acting on a vector from Hilbert space $\mathcal{H}$,
\begin{equation}\label{eq:tdse_it_real}
i\frac{\symrm{d}}{\symrm{d}t}\ket{\Psi(t)}=H\ket{\Psi(t)}.
\end{equation}
Introducing imaginary time $\tau$ via $\tau = i t$ (equivalently $t=-i\tau$) and writing
$\ket{\Psi(\tau)}:=\ket{\Psi(t=-i\tau)}$ gives the imaginary-time evolution equation
\begin{equation}\label{eq:tdse_it}
\frac{\symrm{d}}{\symrm{d}\tau}\ket{\Psi(\tau)}=-H\ket{\Psi(\tau)},
\end{equation}
with formal solution
\begin{equation}\label{eq:wf_it}
\ket{\Psi(\tau)}=e^{-\tau H}\ket{\Psi(0)}.
\end{equation}
Let $\{\ket{\phi_I}\}$ denote a complete orthonormal set of eigenstates of $H$,
\begin{equation}
H\ket{\phi_I}=E_I\ket{\phi_I},
\end{equation}
with real eigenvalues $E_I$ bounded from below. Expanding the initial state in this eigenbasis,
\begin{equation}\label{eq:wf_expansion_eigenstates}
\ket{\Psi(0)}=\sum_I c_I(0)\ket{\phi_I},
\end{equation}
and substituting into Eq.~\eqref{eq:wf_it} yields
\begin{equation}\label{eq:tdse_it_expansion}
\ket{\Psi(\tau)}=\sum_I c_I(0)e^{-\tau E_I}\ket{\phi_I}.
\end{equation}
For $\tau>0$, components with larger eigenvalues $E_I$ decay faster than those with
smaller $E_I$. Assuming the ground-state coefficient $c_0(0)\neq 0$, the normalized state
$\frac{\ket{\Psi(\tau)}}{\lvert\Psi(\tau)\rvert}$ converges to the ground state $\ket{\phi_0}$ in the limit
$\tau\to\infty$. In numerical applications one therefore propagates in small imaginary-time
steps and renormalizes the wavefunction after each step to prevent decay to the zero
vector and to extract the ground state.
